\documentclass[11pt,a4paper]{article}
\usepackage[utf8]{inputenc}
\usepackage[T1]{fontenc}
\usepackage[brazil]{babel}
\usepackage[a4paper,margin=2.3cm]{geometry}
\usepackage{enumitem}
\usepackage{hyperref}
\usepackage{xcolor}
\setlength{\parindent}{0pt}
\setlength{\parskip}{5pt}
\sloppy
\setlength{\emergencystretch}{3em}
\hypersetup{colorlinks=true,urlcolor=blue,linkcolor=black}
\newcommand{\cabecalho}[2]{\begin{center}{\small \textbf{Sprint 2 -- #1}}\\[1mm]\textit{Inventário de Ativos e Vulnerabilidades}\\Matheus Moreli Durão Silva -- 12621CBS217\\Cibersegurança -- Bacharelado -- Noturno\end{center}\hrule\vspace{4mm}}
\newcommand{\bloco}[1]{\par\vspace{2mm}\textbf{#1}\par}
\newcommand{\codigo}[1]{\texttt{#1}}

\begin{document}\cabecalho{S2\_07 -- Especificação Completa do Sistema}{}
\bloco{Escopo} O sistema registra ativos de TI, criticidade, responsável, localização, tipo, vulnerabilidades, fonte, CVE, CWE, CVSS, tratamento e histórico. Permite CRUD de ativos e consultas locais. Não está no escopo: autenticação, varredura automática, integração com APIs externas, interface gráfica ou credenciais.
\bloco{Modelo de dados} \codigo{Ativo}: ID, nome, responsável, localização, tipo, criticidade, vulnerabilidades e histórico. \codigo{Vulnerabilidade}: CVE, CWE, CVSS, descrição, fonte/data, impacto, prioridade, tratamento, status e verificação. Os nomes usados no texto são os mesmos atributos das classes.
\bloco{Entradas e saídas} A entrada acontece no menu do terminal. A saída apresenta confirmação, listagem, resultado de consulta ou mensagem de validação. A persistência é o arquivo \codigo{dados/inventario.json}, gerado sem credenciais.
\bloco{Rastreabilidade} RF01--RF07 e RN01--RN03 estão descritos em S2\_04 e possuem cenários em S2\_06. Assim nenhum requisito fica sem critério de aceitação.
\bloco{Checklist atendido} Escopo declarado, vocabulário coerente, aceitação ligada a cada requisito, campos e regras alinhados e versão preparada para o repositório.
\end{document}
